% 第5章图：消息m_{4->3}的含义。树边{1,3},{2,3},{3,4},{4,5}，根3；删去边{3,4}后4一侧为{4,5}。
% phi4=(0,1), phi5=(3,0)；边{3,4}两端不同成本2，边{4,5}两端不同成本1。
% x3=0：最优(x4,x5)=(0,1)，成本0+0+0+1=1；x3=1：最优(1,1)，成本1+0+0+0=1。故m_{4->3}=(1,1)，回溯b_{4->3}=(0,1)。
\begin{tikzpicture}[x=1mm,y=1mm,>=Stealth,
  v/.style={circle,draw=FigureStroke,fill=FigureBlueFill,line width=.8pt,minimum size=6.5mm,inner sep=0pt},
  w/.style={circle,draw=FigureStroke,fill=white,line width=.8pt,minimum size=6.5mm,inner sep=0pt},
  e/.style={FigureStroke,line width=.8pt},
  lab/.style={font=\figurefont\footnotesize,text=FigureInk},
  note/.style={font=\FigureNoteFont\fontsize{7.5}{9}\selectfont,text=FigureMuted}]
% ---------- tree ----------
\fill[FigureGreenFill,rounded corners=4pt] (21,14) rectangle (47,-14);
\node[font=\FigureNoteFont\fontsize{7.5}{9}\selectfont,text=FigureGreen,anchor=north] at (34,-14.8) {$4$一侧$V_{4|3}$};
\node[w] (t3) at (13,26) {$3$};
\node[w] (t1) at (0,8) {$1$};
\node[w] (t2) at (13,8) {$2$};
\node[v] (t4) at (27,8) {$4$};
\node[v] (t5) at (27,-7) {$5$};
\draw[e] (t1)--(t3) (t2)--(t3) (t4)--(t5);
\draw[e,dashed] (t3)--(t4);
\draw[FigureRed,line width=1pt] ($(t3)!0.5!(t4)+(-2.4,-1.6)$)--($(t3)!0.5!(t4)+(2.4,1.6)$);
\node[note,anchor=south west] at ($(t3)!0.5!(t4)+(1.8,1.2)$) {$2$};
\node[note,anchor=west] at (t4.east) {$\phi_4=(0,1)$};
\node[note,anchor=west] at (t5.east) {$\phi_5=(3,0)$};
\node[note,anchor=east] at ($(t4)!0.5!(t5)+(-1,0)$) {$1$};
% ---------- price list ----------
\begin{scope}[xshift=54mm]
 \node[lab,anchor=west] at (0,27) {消息$m_{4\to3}$：$4$一侧的最低成本};
 \draw[FigureStroke,line width=.6pt] (0,22)--(46,22);
 \node[lab] at (5,18) {$x_3$};
 \node[lab] at (18,18) {最低成本};
 \node[lab] at (36,18) {$(x_4,x_5)$};
 \draw[FigureStroke,line width=.4pt] (0,14.5)--(46,14.5);
 \node[lab] at (5,10) {$0$};
 \node[lab] at (18,10) {$1$};
 \node[lab] at (36,10) {$(0,1)$};
 \node[note] at (18,4.5) {$0+0+0+1$};
 \node[lab] at (5,-2) {$1$};
 \node[lab] at (18,-2) {$1$};
 \node[lab] at (36,-2) {$(1,1)$};
 \node[note] at (18,-7.5) {$1+0+0+0$};
 \draw[FigureStroke,line width=.6pt] (0,-11.5)--(46,-11.5);
\end{scope}
\end{tikzpicture}
